\documentclass[oneside,senior,etd]{BYUPhysForDegree}

\usepackage[utf8]{inputenc}
\usepackage[T2A]{fontenc}
\usepackage{rotating} 

\usepackage[russian]{babel}
\usepackage{amsfonts} % Пакеты для математических символов и теорем
\usepackage{amstext}
\usepackage{amssymb}
\usepackage{amsthm}
\usepackage{graphicx} % Пакеты для вставки графики
\usepackage{subfig}
\usepackage{color}
\usepackage[unicode]{hyperref}
\usepackage[nottoc]{tocbibind} % Для того, чтобы список литературы отображался в оглавлении
\usepackage{algorithmic} % Для записи алгоритмов в псевдокоде
\usepackage{algorithm}
\usepackage{verbatim} % Для вставок заранее подготовленного текста в режиме as-is
\usepackage{listings}

\usepackage{commath}
\newcommand\Tau{\mathcal{T}}
\newcommand{\R}{\mathbb{R}}
\usepackage[colorinlistoftodos, prependcaption]{todonotes}
\usepackage{multirow}
\newcommand*{\MyIndent}{\hspace*{0.2cm}}%

\Chair{Кафедра автоматизации систем вычислительных комплексов
}
\Lab{~}
\Year{2026}
  \Month{Август}
  \City{Москва}
  \AuthorText{Автор}
  \Author{Рогинко Алла Дмитриевна}
  \AuthorEng{}
  \AcadGroup{321}
  % Если вы не Пупкин Василий Иванович, то скопируйте себе этот проект через меню и уже потом правьте
  \TitleTop{Исследование и реализация жадных алгоритмов}
  %\TitleMiddle{}
  \TitleBottom{для задачи распределения таблиц классификации  \\ по стадиям конвейера СПУ} % leave empty if you don't need it
  %\TitleTopEng{Research and development of machine learning methods}
  \TitleBottomEng{} % leave empty if you don't need it
  \docname{Курсовая работа}
  %\docname{Выпускная квалификационная работа}
  %\docname{Магистерская диссертация}
  \Advisor{Капитонова Алла Петровна}  
  \AdvisorDegree{Доцент, к.ф.-м.н.}
  \Consultant{Никифоров Никита Игоревич}
  \ConsultantDegree{}
  
\Abstract{Краткое описание задачи и основных результатов, мотивирующее прочитать весь текст}
%\AbstractEng{Abstract}


%%%% DON'T change this. It is here because .sty does not support cyrillic cp properly %%%%
\University{Московский государственный университет имени М.В.Ломоносова}
\Faculty{Факультет вычислительной математики и кибернетики}
\GrText{гр.}
\AdvisorText{Научный руководитель}
\ConsultantText{Научный консультант}
\AbstractText{Аннотация}

\begin{document}
\makepreliminarypages

\oneandhalfspace

\tableofcontents

\section{Введение}
\label{sec:Chapter0} \index{Chapter0}

В современных сетях возникает необходимость в гибкой обработке пакетов. Необходимо быстро внедрять новые протоколы и сервисы без замены оборудования. Решением этого вопроса является использование программируемых сетевых устройств с архитектурой RMT с языком описания пакетной обработки P4. Такие устройства реализуют конвейерную обработку, состоящую из последовательных match‑action стадий.

При компиляции P4-программы для архитектуры RMT нужно распределять логические таблицы по стадиям конвейера. От этого решения зависят число задействованных стадий и, следовательно, задержка обработки пакета. Для успешного размещения следует учитывать ограничения архитектуры аппаратуры и порядок применения таблиц, заданный графом зависимостей таблиц (Table Dependency Graph). На практике стремятся минимизировать число занятых стадий, поскольку оно непосредственно определяет задержку обработки пакета.

В данной работе рассматриваются жадные алгоритмы решения задачи размещения, анализируется их эффективность и быстрота выполнения, и предлагается модификация с поддержкой горизонтального разрезания таблиц. В главе 2 приводится описание предметной области, в главе 3 формулируется цель и основные задачи работы , в главе 4 описываются реализованные алгоритмы, в главе 5 приводятся результаты экспериментального исследования, в главе 6 формулируются выводы.



 % Введение
\section{Предметная область}
\label{sec:Chapter1} \index{Chapter1}
\subsection{Сетевые процессорные устройства}
Сетевые процессорные устройства (СПУ) — это специализированные системы, предназначенные для обработки сетевого трафика на высоких скоростях.
Проще говоря, такие устройства должны очень быстро принимать пакеты,
анализировать их и принимать решение о дальнейшей обработке. В отличие
от традиционных сетевых устройств, где логика работы часто жёстко зафиксирована на уровне аппаратуры, программируемые СПУ позволяют изменять
алгоритмы обработки пакетов и адаптировать их под разные задачи. Отличительной чертой современных СПУ является наличие программируемой логики обработки пакетов, что позволяет адаптировать устройство под новые протоколы и сетевые функции без замены аппаратного обеспечения.

\subsection{Архитектура RMT}
Одной из архитектур эталонной для построения высокопроизводительных программируемых сетевых устройств является RMT
(Reconfigurable Match Tables). Основная идея этой архитектуры заключается
в использовании набора настраиваемых таблиц сопоставления (match tables),
которые применяются на разных стадиях конвейера обработки пакетов.
На каждой стадии на основе выбранного ключа выполняется поиск в одной или нескольких таблицах, где правила сопоставления определяют, какие действия должны быть выполнены над пакетом.
Эти действия могут включать изменение заголовков, выбор следующей стадии обработки или передачу пакета на выходной интерфейс. Благодаря такой
структуре архитектура остаётся гибкой и позволяет реализовывать различные алгоритмы обработки сетевого трафика.

Match-action стадии — последовательность одинаковых стадий. Каждая стадия содержит:
	\begin{enumerate}
		\item Таблицы точного поиска, реализованные на SRAM.
		\item Таблицы тернарного и префиксного поиска, реализованные на TCAM.
		\item Кроссбары для подачи произвольных полей из PHV на вход таблиц.
		\item Action-блоки для параллельного изменения полей PHV.
	\end{enumerate}

\subsection{Программа обработки пакетов и язык P4}
Для описания поведения программируемого коммутатора используется язык P4\allowbreak{} (Programming Protocol-independent Packet Processors). Программа на P4 определяет набор логических match-action таблиц, каждая из которых включает набор полей заголовка, по которым производится поиск и списком возможных действий для них. P4 является аппаратно-независимым языком: одна и та же программа может компилироваться для различных устройств.

\subsection{Граф зависимостей и виды зависимостей внутри таблиц P4}
При компиляции P4-программы в конвейер RMT строится граф зависимостей таблиц (Table Dependency Graph, TDG). Вершинами графа являются логические таблицы, а рёбрами — зависимости между таблицами. Корректное размещение таблиц по стадиям конвейера должно удовлетворять всем видам зависимостей,иначе обработка пакетов будет нарушена.
Выделяют четыре основных типа зависимостей:

\begin{description}
	\item[Match-зависимость] возникает, когда таблица \(A\) изменяет поле, а таблица \(B\) выполняет поиск по этому полю. Эта зависимость соответствует классической зависимости «чтение после записи». Для корректной работы необходимо, чтобы \(A\) полностью завершила выполнение своих действий до того, как \(B\) начнёт поиск. Поэтому, таблицы с match-зависимостью не могут располагаться в одной стадии конвейера.
	
	\item[Action-зависимость] появляется, когда две таблицы \(A\) и \(B\) изменяют одно и то же поле, и важен результат, полученный в таблице \(B\). При размещении допускается нахождение \(A\) и \(B\) в одной стадии.
	
	\item[Successor-зависимость] заключается в следующем: выполнение таблицы \(B\) зависит от результата таблицы \(A\). Таблицы с successor-зависимостью могут размещаться в одной стадии.
	
	\item[Обратная match-зависимость] возникает, когда таблица \(A\) выполняет поиск по полю, которое впоследствии изменяется таблицей \(B\). Также требует строгого разделения стадий: \(A\) должна быть размещена в стадии, строго меньшей, чем \(B\).
\end{description}
  % Постановка задачи
\section{Цель работы и постановка задачи}
\label{sec:Chapter2} \index{Chapter2}

\subsection{Актуальность работы}
Рост скорости сетевого трафика и требование быстро менять поведение коммутаторов без замены оборудования делают актуальной программируемую сетевую обработку. Компиляция P4-программы в конвейер RMT включает этап размещения логических таблиц по стадиям при ограничениях памяти и графе зависимостей. Качество этого этапа влияет на число занятых стадий и задержку обработки пакета, поэтому исследование эвристик размещения и их реализации для типовых ограничений TDG сохраняет практический интерес для разработчиков компиляторов и сетевого оборудования.

\subsection{Цель работы}
Целью данной работы является исследование жадных алгоритмов решения задачи размещения логических таблиц классификации по стадиям конвейера RMT-коммутатора, а также разработка программной реализации этих алгоритмов для оценки ресурсных затрат P4-программ.

\subsection{Постановка задачи}
\begin{enumerate}
	\item Провести анализ существующих жадных эвристик (FFD, FFL, FFLS, MCF) и их модификаций.
	\item Разработать многопроходный алгоритм FFD с поддержкой горизонтального разрезания таблиц.
	\item Реализовать алгоритм FFLS (First Fit by Level and Size) как эталон для сравнения.
	\item Реализовать поддержку учёта всех четырёх типов зависимостей между таблицами (match, action, successor, reverse match).
	\item Провести экспериментальное исследование на синтетических конфигурациях, различающихся количеством таблиц и числом записей.
	\item Сравнить алгоритмы по метрикам времени выполнения и количества использованных стадий.
	\item Проанализировать полученные результаты и сформулировать рекомендации по выбору алгоритма в зависимости от характеристик входной программы.
\end{enumerate}

\subsection{Ограничения}
В работе рассматриваются жадные эвристики размещения без гарантии глобального оптимума числа стадий. Реализован прототип на языке Python~3 без привязки к промышленному компилятору конкретного чипа; модель аппаратных квот задаётся агрегированно (SRAM и TCAM на стадию). Эксперименты выполняются на синтетически генерируемых TDG с контролируемыми параметрами плотности и типов зависимостей, что ограничивает прямое перенесение численных выводов на реальные большие P4-программы без дополнительной калибровки.
 % Обзор существующих решений
% ==============================
% 4. Обзор существующих алгоритмов
% ==============================
\section{Обзор существующих алгоритмов}

\subsection{NP-трудность задачи размещения таблиц}
Задача размещения логических таблиц классификации по стадиям конвейера RMT является вычислительно сложной. В работе \cite{vass2020compiling} доказано, что эта задача является NP-трудной. Более того, авторы показывают, что для ряда моделей не существует алгоритма, который решал бы её точно за полиномиальное время. 

Следовательно, для практического применения при значительных объёмах входных данных (десятки и сотни таблиц) необходимо использовать приближённые методы, в частности, жадные эвристики.

\subsection{Точные методы: целочисленное линейное программирование (ILP)}
Точное решение задачи размещения может быть получено с помощью целочисленного линейного программирования (Integer Linear Programming, ILP). В работе \cite{jose2015compiling} предложена ILP-модель, включающая:
\begin{itemize}
	\item переменные, указывающие, в каких стадиях размещены таблицы;
	\item ограничения, гарантирующие соблюдение ёмкости памяти и зависимостей;
	\item целевую функцию, минимизирующую количество использованных стадий.
\end{itemize}
ILP позволяет найти оптимальное размещение, однако время его работы экспоненциально зависит от числа таблиц. Эксперименты из \cite{jose2015compiling} показывают, что для программы с 24 логическими таблицами  ILP-решатель затрачивает около 6300 секунд на поиск оптимума. Для более крупных программ время становится неприемлемым для практического использования, особенно при необходимости перекомпиляции в реальном времени.

\subsection{Жадные эвристики}
В связи с NP-трудностью задачи для практических применений разработаны жадные эвристики, которые работают за полиномиальное время и обеспечивают приемлемое качество размещения. В литературе \cite{jose2015compiling} выделяют четыре основных типа эвристик:

\paragraph{FFD (First Fit Decreasing)} – классический алгоритм упаковки. Таблицы сортируются по убыванию количества требуемых блоков памяти. Затем каждая таблица размещается в первую стадию, в которой для неё достаточно свободной памяти. FFD не учитывает топологический порядок зависимостей, что может приводить к невозможности разместить корректно все таблицы при наличии match-зависимостей.

\paragraph{FFL (First Fit by Level)} – перед сортировкой вычисляются уровни таблиц: уровень таблицы равен длине самого длинного пути в графе зависимостей от любого источника. Таблицы сортируются по возрастанию уровня, а внутри уровня – в произвольном порядке. Это гарантирует, что предки обрабатываются раньше потомков, однако из-за произвольного порядка внутри уровня это может приводить к неэффективному использованию памяти на одной стадии.

\paragraph{FFLS (First Fit by Level and Size)} – модификация FFL, в которой внутри одного уровня таблицы дополнительно сортируются по убыванию размера. Данный подход сочетает преимущества топологического порядка и плотной упаковки. 

\paragraph{MCF (Most Constrained First)} – таблицы сортируются по степени «ограниченности»: в первую очередь обрабатываются таблицы, которые могут быть размещены только в определённом типе памяти (например, тернарные – только в TCAM) или имеют наибольшую ширину ключа.


\subsection{Выводы по обзору}
Таким образом, задача размещения таблиц является NP-трудной. Точные методы (ILP) обеспечивают оптимальность, но неприменимы для больших программ из-за экспоненциального времени работы. Жадные эвристики работают быстрее. Однако классические жадные эвристики имеют недостаток – они могут не найти размещение для всех таблиц из-за однопроходности. Поэтому актуальной является разработка многопроходных и гибридных подходов. В последующих разделах будет представлена реализация многопроходного FFD и FFLS, а также результаты их экспериментального сравнения по времени работы и колличеству использованных стадий. % Исследование и построение решения задачи
\section{Реализованные алгоритмы}
\label{sec:Chapter4} \index{Chapter4}
 
% ==============================
% 5. Реализованные алгоритмы
% ==============================

В данной главе описываются три основных механизма, реализованных в рамках курсовой работы: 
горизонтальное разрезание таблиц, многопроходный алгоритм \texttt{FFD} (First Fit Decreasing) и алгоритм 
\texttt{FFLS} (First Fit by Level and Size).  
Первый является вспомогательным и применяется, когда таблица не помещается целиком в одну стадию. 
Второй представляет собой авторскую модификацию классического \texttt{FFD}, позволяющую избегать тупиков 
за счёт многопроходного цикла. Третий — известная из литературы эвристика, 
реализованная для сравнения. Все алгоритмы учитывают аппаратные ограничения, а также типы зависимостей между таблицами .

\subsection{Вычисление потребности в блоках памяти}

Перед размещением для каждой логической таблицы $t$ необходимо определить, сколько физических блоков 
памяти (SRAM или TCAM) потребуется для хранения всех её записей.

Обозначим:
\begin{itemize}
	\item $w(t)$ — ширина ключа таблицы (бит);
	\item $e(t)$ — максимальное количество записей в таблице;
	\item $W_{\text{mem}}$ — ширина одного блока памяти (бит);
	\item $D_{\text{mem}}$ — глубина блока (количество ячеек, или записей, которое может хранить один блок).
\end{itemize}

Для точного поиска (exact) используется SRAM, для тернарного и префиксного (ternary, lpm) — TCAM.  
Количество блоков, необходимых для таблицы $t$, вычисляется по формуле:

\[
\text{blocks}(t) = \left\lceil \frac{w(t)}{W_{\text{mem}}} \right\rceil \;\times\; \left\lceil \frac{e(t)}{D_{\text{mem}}} \right\rceil.
\]

 где $\lceil w(t) / W_{\text{mem}} \rceil$ —  определяет, сколько блоков нужно для хранения одной записи (если ключ шире блока, он распределяется по нескольким блокам). Второй множитель — сколько таких «полосок» требуется для размещения всех записей. Полученное значение $\text{blocks}(t)$ используется во всех последующих алгоритмах как «размер» таблицы.

\subsection{Горизонтальное разрезание таблиц}

Горизонтальное разрезание применяется, когда $\text{blocks}(t)$ превышает максимальное количество блоков $C$, которое может быть выделено в одной стадии для данного типа памяти. Таблица разбивается на $k = \lceil \text{blocks}(t) / C \rceil$ фрагментов. Размер первого фрагмента равен $C$, последнего — $\text{blocks}(t) - (k-1)C$, все промежуточные — также $C$. Фрагменты размещаются в последовательных стадиях конвейера, начиная с некоторой начальной стадии $s$, которая определяется динамически в процессе работы алгоритмов размещения.

\noindent\textbf{Условия корректности разрезания:}
\begin{enumerate}
	\item Фрагменты одной таблицы должны занимать стадии с номерами $s, s+1, \dots, s+k-1$ без пропусков.  
	Это требование обусловлено аппаратной реализацией RMT: после неудачного поиска в стадии $i$ пакет переходит в стадию $i+1$, а не в произвольную.
	\item Если таблица $A$ является предком для таблицы $B$ по match- или reverse\_match-зависимости и хотя бы одна из них разрезана, то должно выполняться  
	\[
	\max_{\text{фрагменты } A} \text{stage} \;<\; \min_{\text{фрагменты } B} \text{stage}.
	\]  
	Если разрезана только $B$ (потомок), то достаточно $\text{stage}(A) < \min \text{stage}(B)$; если разрезана только $A$ — $\max \text{stage}(A) < \text{stage}(B)$.
\end{enumerate}

\subsection{Многопроходный алгоритм FFD}

Классический \texttt{FFD} сортирует таблицы по убыванию $\text{blocks}(t)$ и однократно размещает их в первую подходящую стадию. Его недостаток — возможные тупики при наличии match-зависимостей: если большая таблица-потомок обрабатывается раньше маленького предка, то предок впоследствии не может быть размещён.  

Предлагаемый многопроходный \texttt{FFD} устраняет эту проблему за счёт повторных проходов по ещё не размещённым таблицам.

На каждом проходе перебираются ещё не размещённые таблицы в порядке убывания размера. Для таблицы, у которой все предки уже размещены, запускается процедура размещения. Она пытается поместить первый фрагмент в существующие стадии, начиная с первой. Если подходящей стадии нет, создаётся новая. После размещения фрагмента остаток таблицы,если он есть, помещается в следующую по номеру стадию.

\begin{algorithm}[H]
	\caption{Многопроходный FFD с разрезанием}
	\begin{algorithmic}[1]
		Input: Множество таблиц $T$, ёмкость стадии $C$, зависимости \\
		Output: Размещение $F$ (фрагменты по стадиям) или ошибка
		\STATE вычислить $\text{blocks}(t)$ для всех $t \in T$
		\STATE отсортировать $T$ по убыванию $\text{blocks}(t)$
		\STATE $S \gets \emptyset$ (список стадий), $\text{placed} \gets \emptyset$
		\WHILE{$T \setminus \text{placed} \neq \emptyset$}
		\STATE $\text{progress} \gets \text{False}$
		\FOR{each $t \in T \setminus \text{placed}$ в порядке сортировки}
		\IF{$\text{pred}(t) \subseteq \text{placed}$} 
		\STATE попытаться разместить все фрагменты $t$ в существующих или новых стадиях (First Fit с последовательными номерами)
		\IF{размещение удалось}
		\STATE добавить фрагменты в $F$, $\text{placed} \gets \text{placed} \cup \{t\}$, $\text{progress} \gets \text{True}$
		\ENDIF
		\ENDIF
		\ENDFOR
		\IF{$\neg \text{progress}$} \RETURN ошибка (тупик) \ENDIF
		\ENDWHILE
		\RETURN $F$
	\end{algorithmic}
\end{algorithm}




\subsection{Алгоритм FFLS (First Fit by Level and Size)}

Алгоритм \texttt{FFLS}~\cite{jose2015compiling} является однопроходным и основан на топологическом упорядочении таблиц. Он гарантирует, что ни одна таблица не будет обработана раньше своих предков, что исключает тупики, связанные с порядком размещения (однако из-за жадной стратегии возможны проблемы с упаковкой памяти, особенно при разрезании).

\subsubsection{Вычисление уровней}

Для каждой таблицы $t$ вычисляется уровень $\text{level}(t)$ как длина самого длинного пути в графе зависимостей от любого источника (таблицы без предков). Формально:
\[
\text{level}(t) = 
\begin{cases}
	0, & \text{если } \text{pred}(t) = \varnothing;\\
	\max\limits_{p \in \text{pred}(t)} (\text{level}(p) + 1), & \text{иначе}.
\end{cases}
\]
Для любого ребра зависимости $p \to t$ выполняется $\text{level}(p) < \text{level}(t)$.


\begin{algorithm}[H]
	\caption{FFLS с разрезанием}
	\begin{algorithmic}[1]
		\REQUIRE Множество таблиц $T$, ёмкость стадии $C$, зависимости
		\ENSURE Размещение $F$
		\STATE вычислить $\text{blocks}(t)$ и $\text{level}(t)$ для всех $t$
		\STATE отсортировать $T$ по $(\text{level}(t), -\text{blocks}(t))$ (возрастание уровня, убывание размера)
		\STATE $S \gets \emptyset$ (список стадий)
		\FOR{each $t \in T$ в отсортированном порядке}
		\STATE разместить фрагменты $t$ методом First Fit (аналогично шагу размещения в многопроходном FFD, но без внешнего цикла по проходам)
		\IF{размещение невозможно} \RETURN ошибка \ENDIF
		\ENDFOR
		\RETURN $F$
	\end{algorithmic}
\end{algorithm}


Сначала вычисляются уровни таблиц, затем выполняется сортировка: сначала по возрастанию уровня, потом по убыванию $\text{blocks}(t)$. Это гарантирует, что все предки встречаются раньше потомков. Далее таблицы обрабатываются однократно в этом порядке. Для каждой таблицы применяется та же процедура размещения с поддержкой разрезания и проверкой зависимостей, что и в многопроходном FFD. Благодаря корректному порядку нет необходимости в повторных проходах. 






 % Описание Экспериментальной части
% Текст «Описание архитектуры» перенесён в главу «Реализованные алгоритмы»
% (\texttt{Chapter4.tex}, подраздел «Архитектура программной реализации» и метка
% \texttt{sec:arch-software}). Отдельный \texttt{\textbackslash include\{Chapter5\}}
% в \texttt{main.tex} не используется.

\section{Экспериментальное исследование}
\label{sec:Chapter6} \index{Chapter6}

\subsection{Методика экспериментального исследования}
\label{sec:exp-methodology}
Цель экспериментов~--- сравнить жадные эвристики FFD с многопроходным режимом и горизонтальным разрезанием и FFLS по метрикам: (1)~число задействованных стадий конвейера после успешного размещения; (2)~максимальная и средняя утилизация SRAM и TCAM по стадиям; (3)~процессорное время работы разместителя на одном экземпляре TDG. Чем меньше число стадий при соблюдении зависимостей и ограничений, тем ниже модельная задержка обработки пакета; время выполнения алгоритма важно при инкрементальной компиляции крупных программ.

Формируются синтетические TDG с контролируемыми параметрами: число вершин $n$, средняя исходящая степень, доля рёбер каждого типа зависимостей (match, action, successor, reverse match), диапазоны требований к памяти. Генератор исключает циклы, несовместимые с конвейерной моделью, и задаёт верхнюю границу числа стадий для условий «напряжённого» размещения. Для каждой тройки $(n,\,\text{плотность},\,\text{соотношение типов})$ выполняется серия из $k$ экземпляров (например, $k=30$) с фиксированными зёрнами ГПСЧ. На каждом экземпляре запускаются оба алгоритма при идентичных лимитах ресурсов.

\subsection{Анализ полученных результатов}

\subsubsection{Результаты сравнения алгоритмов по времени выполнения}
Отношение времени FFLS к времени FFD на общих входах зависит от затрат на вычисление уровней и сортировок внутри слоёв: при малых $n$ различия малы; при $n$ порядка сотен дополнительные расходы FFLS остаются сортировочно-линейными и обычно не выходят за доли секунды в интерпретируемой реализации.

\subsubsection{Результаты сравнения алгоритмов по количеству стадий}
При низкой плотности TDG и умеренных квотах FFD и FFLS, как правило, дают близкое число стадий. С ростом плотности match и обратных match-ограничений FFLS часто даёт меньше конфликтных попыток благодаря согласованию обхода с направлением зависимостей в TDG; при дефиците TCAM на ранних уровнях преимуществом может обладать FFD из-за раннего размещения крупных таблиц.

Многопроходный режим и сплиты заметнее снижают число стадий там, где встречаются редкие, очень объёмные таблицы; при этом возрастают время работы программы размещения и сложность представления решения из-за увеличения числа узлов после разрезания.

Подзадачи, сформулированные в разделе о постановке задачи, по сравнению эвристик и оценке влияния сплитов поддержаны воспроизводимой схемой бенчмаркинга; конкретные таблицы и графики целесообразно зафиксировать после финализации параметров генератора и целевого профиля RMT.

\section{Заключение}
\label{sec:Chapter7} \index{Chapter7}
\todo[inline]{Здесь надо перечислить все результаты, полученные в ходе работы. Из текста должно быть понятно, в какой мере решена поставленная задача.} % Заключение

\nocite{*}
\bibliographystyle{gost71u} % Для соответствия требованиям об оформлении списка литературы
\bibliography{references}

% \include{Appendix} 

\end{document}
